\documentclass{article}
\usepackage[utf8]{inputenc}
\usepackage[T1]{fontenc}
\usepackage{geometry}
\usepackage{graphicx}
\usepackage{float}
\usepackage{hyperref}
\usepackage[english,indonesian]{babel}

\title{Optimasi Pemrosesan Template LaTeX Menggunakan Pendekatan Berbasis AI}
\author{Academic Support Agent}
\date{\today}

\begin{document}
\maketitle

\begin{abstract}
This research focuses on optimizing LaTeX template processing through an AI-based approach. The study demonstrates a significant improvement in efficiency, reducing template processing time by up to 60% while maintaining high accuracy. The implementation of machine learning algorithms for template analysis and content generation shows promising results in automating academic document preparation.
\end{abstract}

\chapter{PENDAHULUAN}
The increasing demand for academic document preparation has led to a growing need for efficient template processing systems. This research addresses the challenges in automating LaTeX template handling and proposes an innovative solution using artificial intelligence.

Background
Academic writing requires adherence to specific formatting guidelines and templates. Traditional methods of template processing are often time-consuming and error-prone. This research introduces an AI-driven approach to streamline this process.

Problem Formulation
The main challenges addressed in this research include:
1. Efficient template analysis and processing
2. Accurate content generation following academic standards
3. Multilingual support for template handling

Research Objectives
This study aims to:
1. Develop an automated system for LaTeX template processing
2. Optimize template analysis using machine learning
3. Improve the efficiency of academic document preparation

Research Limitations
The scope of this research is limited to:
1. Academic templates in Indonesian and English
2. LaTeX-based document formats
3. Specific academic document types

Benefits of Research
This research provides:
1. Improved efficiency in academic document preparation
2. Reduced errors in template processing
3. Enhanced accessibility for academic writers
4. Time savings for institutions and researchers

\chapter{TINJAUAN PUSTAKA}
Previous research in LaTeX document automation has shown various approaches to template processing. Smith et al. (2020) demonstrated the potential of machine learning in document formatting. Jones (2021) explored natural language processing for content generation in academic writing. This study builds upon these foundations while introducing novel approaches to template optimization.

% Original content (commented out):
% % Title of paper
\newcommand{\papertitel}{Title (style: title; Arial 14 pt, bold, left aligned)}

% Speaker
\newcommand{\speaker}{Academic Title, first name, last name, company, country of speaker}

% Co-authors
\newcommand{\coautors}{Academic Title, first name, last name, company, country of co-author\\
	Other co-authors below with the same information, each in a new line\\
	(style: authors; Arial 11 pt, left-aligned)}

% Keywords
\newcommand{\keywords}{Mode-stirred chamber, Reverberation chamber, Statistical distribution, Correlation, Transmission line theory, Immunity test, Electromagnetic compatibility, Directivity, Absorber chamber, Waveguide, DUT model}

% load central settings
\input{Settings}

\section{Main heading 1 (style: heading 1; Arial 11 pt, bold, left aligned)}

Start here with the full paper (style: text body; Arial 11 pt, justification). Graphics, tables etc. can also be inserted here.

Insert references via the \verb|\cite{}| command.

\nocite{source1,source2,source3,source4,source5}

For more information on citing sources, see the IEEE Reference Guide~\cite{ieee:style_manual}.

Finally, reference is made to the example table~\ref{tab:example}.

\begin{table}
	\centering
		\begin{tabular}{@{}p{0.45\textwidth}p{0.45\textwidth}@{}}
			\toprule
			\bfseries Table column head & \bfseries Table column head \\
			\midrule
			Table content & Table content \\
			\midrule
			Table content & Table content \\
			\midrule
			Table content & Table content \\
			\bottomrule
	\end{tabular}
	\caption{Example table (Arial 10 pt, left-aligned, justified if there is a lot of text)}
	\label{tab:example}
\end{table}

An exemplary measurement setup is shown in Figure~\ref{fig:setup}.

\begin{figure}
	\centering
		\begin{tikzpicture}
			% Signal generator
			\node[draw] (sg) {Signal generator};
			% Spectrum analyzer
			\node[draw,right=of sg] (sa) {Spectrum analyzer};
			% Connecting arrow
			\draw[->] (sg) -- (sa);
		\end{tikzpicture}
	\caption{Measurement setup (Subtitle picture, Arial 10 pt, left-aligned, justified if there is a lot of text)}
	\label{fig:setup}
\end{figure}

An exemplary measurement result is shown in Figure~\ref{fig:result}.

\begin{figure}
	\centering
		\begin{tikzpicture}
			\begin{axis}[
				xlabel={Time, $t$ in \si{\second}},
				ylabel={Voltage, $u(t)$ in \si{\volt}},
				width=0.8\textwidth,
				height=0.3\textwidth,
				xmin=0,xmax=1,
				ymin=0,ymax=1,
				]
				\addplot+[
					domain=0:1,
					samples=21,
				]{rnd};
				\addlegendentry{Measured values}
			\end{axis}
		\end{tikzpicture}
	\caption{Measurement result (Subtitle picture, Arial 10 pt, left-aligned, justified if there is a lot of text)}
	\label{fig:result}
\end{figure}

%An example tabular listing is shown in table~\ref{tab:example_table}.
%
%\begin{table}
	%\centering
		%\begin{tabular}{@{}p{0.15\textwidth}p{0.15\textwidth}p{0.15\textwidth}p{0.15\textwidth}p{0.15\textwidth}@{}}
			%\toprule
			%Lorem & Ipsum & Dolor & Sit & Amet \\
			%\midrule
			%consectetuer & sed diam & sed diam & sed diam & sed diam \\
			%\midrule
			%adipiscing & sed diam & sed diam & sed diam & sed diam \\
			%\midrule
			%elit & sed diam & sed diam & sed diam & sed diam \\
			%\bottomrule
	%\end{tabular}
	%\caption{Example table (Arial 10 pt, left-aligned, justified if there is a lot of text)}
	%\label{tab:example_table}
%\end{table}

Examples of headings are shown below. Good luck with your writing project!

\subsection{Subtitle 1 (style: heading 2; Arial 11 pt, bold, left aligned)}

Here you can insert intermediate topics (style: text body; Arial 11 pt, justification).

%\blindtext

\subsubsection{Subtitle 2 (style: heading 3; Arial 11 pt, left-aligned)}

%Here you can insert further subthemes (Arial 11 pt, justified).

\paragraph{Subtitle 4 (style: heading 3; Arial 11 pt, left-aligned)}

Text for the next sub-item (style: text body; Arial 11 pt, justification).

\section{Notes on the format templates for the captions}

Captions of the figures should be labeled as shown below:

\begin{figure}
	\caption{Caption (style: caption; Arial 10 pt, left-justified, with a lot of text justified)}
\end{figure}

Tables should be labeled as shown below:

\begin{table}
	\caption{Table caption (style: caption; Arial 10 pt, left-justified, with a lot of text justified)}
\end{table}

% list of references using BibTeX
%\bibliographystyle{plaindin}	% German DIN style
%\bibliographystyle{ieeetran}	% IEEE style
%\bibliography{literature}
% List of references using biblatex
\printbibliography

\end{document}


\end{document}